% Generated draft: requires mathematical and source review.
\documentclass[11pt]{article}
\usepackage[margin=25mm]{geometry}
\usepackage{fontspec}
\setmainfont{Latin Modern Roman}
\usepackage{amsmath,amssymb,mathtools,unicode-math}
\setmathfont{Latin Modern Math}
\usepackage{enumitem,needspace,xurl,hyperref}
\setlength{\parindent}{0pt}
\setlength{\parskip}{5pt}
\setlength{\emergencystretch}{3em}
\newcommand{\N}{\mathbb{N}}
\newcommand{\Z}{\mathbb{Z}}
\newcommand{\Q}{\mathbb{Q}}
\newcommand{\R}{\mathbb{R}}
\newcommand{\C}{\mathbb{C}}
\newcommand{\F}{\mathbb{F}}
\newcommand{\A}{\mathbb{A}}
\newcommand{\PP}{\mathbb{P}}
\newcommand{\E}{\mathbb{E}}
\newcommand{\eps}{\varepsilon}
\newcommand{\dd}{\,\mathrm{d}}
\begin{document}
\section*{UnsolvedMath ID: 30003489}
\textbf{Source problem code:} catalogue entry
\subsection{Definitions and mathematical statement}
The supplied record does not specify the tensor dimensions, the observed entries, or the precise hypersurface \(H\). The following gives a conditional formulation using the usual finite discrete probability model; identifying it with the source's intended model requires review.

Let \(d\geq 1\), let \(n_1,\ldots,n_d\geq 1\), and write \([n]=\{1,\ldots,n\}\). For each \(k\), define the probability simplex
\[
\Delta_{n_k-1}=\left\{u\in\R_{\geq0}^{n_k}:\sum_{j=1}^{n_k}u_j=1\right\}.
\]
A joint distribution of independent random variables taking values in \([n_1],\ldots,[n_d]\) is a tensor \(p\) with entries
\[
p_{i_1,\ldots,i_d}=\prod_{k=1}^{d}u^{(k)}_{i_k},\qquad u^{(k)}\in\Delta_{n_k-1}.
\]
For an observed index set \(E\subseteq\prod_{k=1}^{d}[n_k]\), a nonnegative partial tensor is a vector \(x\in\R_{\geq0}^{E}\). It is completable if there exist such probability vectors satisfying
\[
x_{i_1,\ldots,i_d}=\prod_{k=1}^{d}u^{(k)}_{i_k}\qquad\text{for every }(i_1,\ldots,i_d)\in E.
\]
Denote the set of completable partial tensors by \(C_E\).

One possible convention for its algebraic boundary is the complex Zariski closure
\[
B_E=\overline{\partial_{A_E}C_E}^{\,\mathrm{Zar}}\subseteq (A_E)_{\C},\qquad A_E=\operatorname{aff}_{\R}(C_E),
\]
where \(\partial_{A_E}\) is the topological boundary relative to the real affine hull and \((A_E)_{\C}\) is its complexification. If the intended \(H\) is a reduced hypersurface in this ambient space, its degree means the total degree of a square-free defining polynomial, equivalently the degree of its projective closure.

\textbf{Degree question.} For the shapes and observed index sets intended by the source, determine \(\deg H\), where \(H\) is the specified algebraic-boundary hypersurface of \(C_E\). The record does not identify these shapes and index sets, establish that \(B_E\) is a hypersurface, or specify whether \(H=B_E\) or only a selected boundary component. Consequently, it does not determine a unique hypersurface or degree problem.

\textbf{Additional question in the original statement.} Is the set
\[
N_E=\R_{\geq0}^{E}\setminus C_E
\]
always convex? Under a universal interpretation over the finite discrete models above, this asks whether, for every \(d,n_1,\ldots,n_d,E\), every \(x,y\in N_E\), and every \(t\in[0,1]\), one has \(tx+(1-t)y\in N_E\). The source record does not specify the intended range of models or whether a smaller ambient domain of admissible partial tensors is imposed. This convexity question occurs in the original statement but is omitted from the clean statement; its inclusion in the intended target remains unresolved.
\subsection{Short English statement}
Given selected entries of a joint probability table, ask whether they extend to a distribution of independent random variables. The main question asks for the degree of the hypersurface describing the algebraic boundary of the extendable entries. The original record also asks whether the nonnegative entries that cannot be extended always form a convex set. Tensor sizes, observation patterns, the precise boundary hypersurface, and the scope of the convexity question require clarification.
\subsection{Sources}
\begin{itemize}
\item Supplied local UnsolvedMath registry record, permanent ID \(30003489\), source problem code \texttt{catalogue entry}, titled ``Degree of the Algebraic Boundary in Partial Tensor Completion.'' The record gives proposed year 2017 and names no proposer.
\item Supplied source citation: \emph{Algebraic Statistics (2018)}, Oberwolfach Reports, pp.\ 1207--1279, DOI \texttt{10.4171/owr/2017/20}. \url{https://doi.org/10.4171/owr/2017/20}. The background describes Oberwolfach Reports as documenting mathematical workshops and their open-problem discussions. The literature-source label is ``Oberwolfach Reports 2018.'' The supplied 2018 report label and 2017 DOI/year metadata are retained without reconciliation.
\item Supplied literature source: ``Kahle et al., rank-one tensor completion.'' \url{https://doi.org/10.1137/16M1074102}. No complete author list, precise title, or relevant theorem location is supplied.
\item Literature checked at: \texttt{2026-08-21}. This is supplied dated curation, not an independently verified assessment. The record labels the problem \emph{partially solved} and states that the algebraic boundary and several diagonal or low-dimensional cases are known, while the general degree question and universal convexity question remain unresolved. Its background explicitly requests independent verification of time-sensitive status and cited claims.
\item The registry labels the statement \texttt{corrected\_verified} and reports that the official Oberwolfach report was checked directly. These are supplied provenance claims; no independent source visit or verification is asserted here.
\end{itemize}
\end{document}
